% !TeX root = ../main.tex
% Research framework for Part 2, redrawn 2026-10-29 from the confirmed v3 preview.
% Layout: overall topic, two narrow scientific-question cards over the regions
% they govern, and four content columns with one problem box each. The box in a
% column states the single problem that this content is to solve; it is not a
% task list. No numbers, arrows, data flows or validation blocks.
\begingroup
\linespread{1}\selectfont
\definecolor{rfDark}{HTML}{2F5A86}
\definecolor{rfLine}{HTML}{4E779E}
\definecolor{rfTint}{HTML}{E4EEF7}
\begin{tikzpicture}[
  x=1cm,y=1cm,line cap=round,line join=round,text=black,
  font=\songti\fontsize{9pt}{11pt}\selectfont,
  every node/.style={align=center,outer sep=0pt,inner sep=0pt},
  rfTitle/.style={font=\heiti\fontsize{12.4pt}{14.8pt}\selectfont,text=white},
  rfQTag/.style={font=\heiti\fontsize{9pt}{11pt}\selectfont,text=rfDark},
  rfQ/.style={font=\songti\fontsize{9pt}{10.8pt}\selectfont},
  rfHead/.style={font=\heiti\fontsize{8.6pt}{10.2pt}\selectfont,text=white},
  rfProb/.style={font=\heiti\fontsize{8pt}{9.6pt}\selectfont,text=rfDark},
  rfProbText/.style={font=\songti\fontsize{8.4pt}{10.2pt}\selectfont},
  rfBracket/.style={font=\songti\fontsize{8.6pt}{10pt}\selectfont}
]
  \path[use as bounding box] (0,3.50) rectangle (16.10,10.92);

  % 1. Overall topic of the proposal.
  \path[fill=rfDark] (0.10,10.08) rectangle (16.00,10.86);
  \node[rfTitle] at (8.05,10.47) {连续动态世界模型构建方法研究};

  % 2. Two narrow question cards, each placed over the region it governs: the
  %    first centred above contents 1--3, the second above content 4. Both use
  %    the same frame height and type size; a short stem and a distribution bar
  %    then mark the mapping (no arrows, no flow labels).
  \path[fill=rfTint,draw=rfLine,line width=.8pt] (2.20,8.40) rectangle (9.80,9.95);
  \node[rfQTag] at (6.00625,9.66) {科学问题一};
  \node[rfQ] at (6.00625,8.92) {不完备观测下连续运动与场景结构的\\可区分条件};

  \path[fill=rfTint,draw=rfLine,line width=.8pt] (12.00,8.40) rectangle (16.00,9.95);
  \node[rfQTag] at (14.00,9.66) {科学问题二};
  \node[rfQ] at (14.00,8.86) {动作响应差异的表征依赖\\与误差传递机理};

  % Mapping: the symmetric bar below the first card feeds contents 1--3; the
  % second card drops straight onto content 4.
  \draw[rfLine,line width=.9pt] (6.00625,8.40) -- (6.00625,8.20);
  \draw[rfLine,line width=.9pt] (1.91875,8.20) -- (10.09375,8.20);
  \draw[rfLine,line width=.9pt] (1.91875,8.20) -- (1.91875,7.90);
  \draw[rfLine,line width=.9pt] (6.00625,8.20) -- (6.00625,7.90);
  \draw[rfLine,line width=.9pt] (10.09375,8.20) -- (10.09375,7.90);
  \draw[rfLine,line width=.9pt] (14.18125,8.40) -- (14.18125,7.90);

  % 3. Four content columns: a tinted panel, a dark header and one problem box.
  %    The box states the single problem of that content; the four names are the
  %    four subsections of the text.
  \foreach \xl/\xr in {%
    0.1000/3.7375,%
    4.1875/7.8250,%
    8.2750/11.9125,%
    12.3625/16.0000}{
    \path[fill=rfTint] (\xl,5.14) rectangle (\xr,7.90);
    \path[fill=rfDark] (\xl,7.12) rectangle (\xr,7.90);
  }

  \node[rfHead] at (1.91875,7.51) {不确定性感知的\\连续时间场景表征};
  \node[rfHead] at (6.00625,7.51) {连续运动引导的\\空间结构表征};
  \node[rfHead] at (10.09375,7.51) {物理连接变化下的\\连续动态场景表征};
  \node[rfHead] at (14.18125,7.51) {基于连续动态表征的\\动作后果预测};

  % One problem box per column.
  \foreach \xl/\xr in {0.1900/3.6475,4.2775/7.7350,8.3650/11.8225,12.4525/15.9100}{%
    \path[fill=white,draw=rfLine,line width=.5pt] (\xl,7.04) rectangle (\xr,5.24);
  }

  \node[rfProb] at (1.91875,6.86) {研究问题};
  \node[rfProbText] at (1.91875,6.03)
    {有限观测与运动歧义下\\如何恢复非采样时刻的\\运动、几何及其不确定性};

  \node[rfProb] at (6.00625,6.86) {研究问题};
  \node[rfProbText] at (6.00625,6.03)
    {遮挡与共动造成归属歧义\\如何确定对象或部件归属\\并限定信息共享范围};

  \node[rfProb] at (10.09375,6.86) {研究问题};
  \node[rfProbText] at (10.09375,6.03)
    {切割撕裂后原有联系失效\\如何延续稳定区域历史\\并表达新生边界};

  \node[rfProb] at (14.18125,6.86) {研究问题};
  \node[rfProbText] at (14.18125,6.03)
    {相近几何下动作后果不同\\如何由运动与当前结构\\约束给定动作的演化预测};

  % 4. Two brackets: the first three contents build one representation, and
  %    that representation stays connected with the prediction content.
  \draw[rfLine,line width=.9pt]
    (1.91875,5.14) -- (1.91875,4.88) -- (10.09375,4.88) -- (10.09375,5.14);
  \node[rfBracket] at (6.00625,4.62) {连续动态表征的协同构建};
  \draw[rfLine,line width=.9pt]
    (6.00625,4.40) -- (6.00625,4.04) -- (14.18125,4.04) -- (14.18125,5.14);
  \node[rfBracket] at (10.09375,3.76) {表征与预测的内在联系};
\end{tikzpicture}
\endgroup
